\section{Evaluation}
\label{sec:eval}

We evaluate \sys along five questions:

\begin{itemize}[leftmargin=*]
  \item \textbf{RQ1:} What is the runtime overhead of the \sys runtime compared
        to existing GPU instrumentation frameworks?
  \item \textbf{RQ2:} How efficiently can \sys maintain shared state across the
        CPU–GPU boundary?
  \item \textbf{RQ3:} Can \sys express practical GPU observability tools with
        significantly lower overhead than CUPTI and NVBit?
  \item \textbf{RQ4:} Can \sys-based memory policies improve performance over
        default UVM behavior across different access patterns?
  \item \textbf{RQ5:} Can \sys-based scheduling policies provide tangible gains
        in tail latency, throughput, and power efficiency in mixed workloads?
\end{itemize}

\subsection{Methodology and Setup}

We evaluated \sys on two machines. Server~A is a single-socket Intel Core Ultra~9
285K processor (24 cores) with 128~GB DDR5 RAM and an NVIDIA RTX~5090 GPU.
Server~B is a single-socket Intel Xeon~6787P (86~cores, 2.0~GHz, 336~MB LLC)
with 768~GB DDR5 RAM (512~GB CXL module) and an NVIDIA H100 GPU. Unless
otherwise noted, we report results on Server~A and average each metric over
10 trials; when averaging speedups we use geometric means. We compare against
NVIDIA's CUPTI and NVTX-based profiler stack, NVBit gpumemtrace, default UVM
memory management, and the native GPU scheduler without \sys-based policies.

\subsection{RQ1: Runtime Overhead vs Existing Tools}

We first quantify the overhead of the \sys runtime itself, independent of any
specific policy logic. We use two microbenchmarks:

\begin{itemize}[leftmargin=*]
  \item A load/store latency benchmark that issues controlled memory accesses
        from GPU kernels, instrumented with \sys and with NVBit gpumemtrace.
  \item A ResNet inference workload instrumented with \sys-based tracing and
        with the PyTorch profiler (CUPTI+NVTX).
\end{itemize}

\Cref{fig:runtime-overhead} shows the per-access latency slowdown under
\sys and NVBit as access sizes increase. \sys adds substantially less latency
than NVBit, remaining low and stable below 128~KB and increasing modestly
thereafter. \Cref{fig:resnet-latency} compares end-to-end ResNet inference
latency distributions: p10 and mean latencies are slightly better than under the PyTorch profiler, and p99 latency is comparable. In both
cases, lightweight launch-time injection and warp-level execution of \sys
handlers are key to keeping overhead low.

\begin{figure} \centering \fbox{\parbox{0.5\textwidth}{\centering [Figure: Runtime Overhead]}} \caption{Runtime Overhead} \label{fig:runtime-overhead} \end{figure}

\begin{figure} \centering \fbox{\parbox{0.5\textwidth}{\centering [Figure: Resnet Inference Latency Overhead]}} \caption{Resnet Inference Latency Overhead} \label{fig:resnet-latency} \end{figure}

\subsection{RQ2: Shared State and Cross-domain Overheads}

Next, we measure the cost of maintaining shared state between CPU and GPU via
\sys maps. We benchmark shared-memory bandwidth and one-way latency on both
servers using a 1~GiB region and both GPU-to-CPU and CPU-to-GPU access patterns.
Servers~A and~B deliver 12–18~GB/s over PCIe with 8–15~µs cross-domain latencies,
which is sufficient for the policy state sizes we maintain. Our SIMT-aware
aggregation and sampling strategies further reduce the frequency of cross-domain
synchronization, as we validate in the use cases below.

\subsection{RQ3: Low-overhead GPU Observability}

We evaluate the suite of bcc-style observability tools built on \sys, measuring
accuracy, overhead, and practical utility across diverse workloads.

\paragraph{Thread divergence and workload balance.}
\texttt{kernelretsnoop} attaches to kernel return paths and records per-thread
timestamps, revealing divergence. Applied to a matrix multiply kernel with
boundary checks, it shows that threads handling edge cases (5% of total) incur
600–850~ns longer execution than interior threads, with timestamps within 10~ns
of hardware performance counters and 2.1% overhead. Refactoring boundary
handling based on this insight improves overall kernel performance by 8%.
\texttt{threadhist} reveals load imbalance in a particle simulation, identifying
a bimodal distribution of iterations per thread that, when corrected, yields an
18% throughput improvement with 1.8% tool overhead.

\paragraph{Launch latency and memory behavior.}
\texttt{launchlate} combines host and device hooks to measure kernel launch
latency. On a pipeline of 50 small kernels, default synchronous execution shows
150–550~µs launch latency per kernel, leading to 32~ms total runtime; switching
to CUDA graphs guided by \texttt{launchlate} reduces this to 7.2~ms. The tool
adds less than 0.5% overhead. \texttt{mem_trace} instruments sampled memory
operations and recovers per-warp bank conflict rates within 2% of hardware
counters, with under 3% end-to-end overhead across GEMM, SpMV, and stencil
kernels (\Cref{fig:obs-overhead}).

\begin{figure}
\centering
\fbox{\parbox{\columnwidth}{\centering [Figure: GPU Observability Tools Overhead Comparison]}}
\caption{Overhead comparison for bcc-style GPU observability tools across different workloads}\label{fig:obs-overhead}
\end{figure}


\paragraph{Comparison to CUPTI and NVBit.}
\Cref{fig:obs-comparison} compares \sys-based tools against CUPTI and
NVBit for common profiling tasks. CUPTI-based profiling incurs 15–25% overhead
and typically requires kernel restarts; NVBit adds 30–45% overhead for similar
metrics. \sys tools average 2–3% overhead and support live instrumentation
without interrupting execution.

\begin{figure}
\centering
\fbox{\parbox{\columnwidth}{\centering [Figure: Observability Tool Comparison]}}
\caption{Comparison of \sys observability tools vs. CUPTI and NVBit for common profiling tasks}\label{fig:obs-comparison}
\end{figure}


\subsection{RQ4: Programmable Memory Management}

We evaluate \sys's memory policy interface on workloads with diverse access
patterns: sequential (vector addition), strided (matrix transpose), and random
(graph traversal). We compare three configurations: (1) default NVIDIA UVM with
on-demand page migration, (2) static prefetching with manual hints, and (3) \sys
with adaptive region-based prefetch and eviction policies that use both host-side
events and device-side memory events.

\Cref{fig:uvm-perf} shows normalized execution time across workloads.
For sequential access, \sys achieves up to 1.8$\times$ speedup over default UVM
by prefetching consecutive regions based on observed access and device hints.
For strided patterns, \sys learns strides and prefetches accordingly, achieving
1.4$\times$ speedup. Random access sees a modest 1.2$\times$ improvement from
working-set tracking and retention of hot regions.

\Cref{fig:page-faults} shows page fault rates. Default UVM incurs 15–30K
faults per second depending on the workload; \sys reduces this to 5–12K~faults/s
(50–65% reduction) through predictive prefetching and more selective eviction.
The adaptive policy adds only 2–3% instrumentation overhead while eliminating
the majority of fault-induced stalls. Average memory footprint increases by
10–15%, which we bound using a hybrid policy that throttles prefetch aggressiveness
under tight capacity constraints, keeping performance within 5% of an
unbounded-memory configuration.

\begin{figure}
\centering
\fbox{\parbox{\columnwidth}{\centering [Figure: UVM Performance Comparison]}}
\caption{Normalized execution time for UVM management strategies across different access patterns}\label{fig:uvm-perf}
\end{figure}

\begin{figure}
\centering
\fbox{\parbox{\columnwidth}{\centering [Figure: Page Fault Frequency]}}
\caption{Page fault frequency comparison between default UVM and \sys programmable UVM}\label{fig:page-faults}
\end{figure}

\subsection{RQ5: Fine-grain Scheduling and Power-aware Control}

Finally, we evaluate \sys's scheduling interface using two classes of policies:
priority-based preemption and power-aware throttling, on mixed workloads that
combine high-priority latency-sensitive kernels (real-time video processing)
with low-priority batch ML training jobs.

\paragraph{Priority-based preemption.}
We compare three approaches: (1) no prioritization (FIFO), (2) CUDA stream-level
priorities, and (3) \sys-based warp-level preemption. \Cref{fig:p99-latency}
plots p99 latency for high-priority tasks. FIFO exhibits 25–80~ms p99 latency
due to blocking behind long-running kernels. CUDA streams reduce p99 to 15–30~ms
but still suffer head-of-line blocking. \sys cuts p99 latency to 5–12~ms by
preempting low-priority warps at synchronization points and handing resources
to high-priority work, with 4–6% overhead and 0.3–0.8~µs per scheduling decision.
Low-priority throughput drops by 8–12%, but \sys maintains better fairness than
CUDA streams, which incur 18–25% throughput reduction due to prolonged starvation
(\Cref{fig:lowprio-throughput}).

\begin{figure}
\centering
\fbox{\parbox{\columnwidth}{\centering [Figure: Priority Scheduler Tail Latency]}}
\caption{P99 latency for high-priority tasks under different scheduling policies}\label{fig:p99-latency}
\end{figure}

\begin{figure}
\centering
\fbox{\parbox{\columnwidth}{\centering [Figure: Low-Priority Task Throughput]}}
\caption{Throughput impact on low-priority tasks with different scheduling policies}\label{fig:lowprio-throughput}
\end{figure}

\paragraph{Power-aware throttling.}
We implement a thermal-aware scheduler that throttles thread blocks when GPU
temperature exceeds configured thresholds. \Cref{fig:power-throttle} shows
power consumption and performance over time under sustained load. Without
throttling, power spikes to 320~W and violates thermal limits. NVIDIA's frequency
scaling reduces power to 280~W but lowers performance by 20%. \sys's selective
thread throttling maintains 290~W average power while reducing performance by
only 8%, improving performance-per-watt compared to frequency scaling. The
sub-microsecond decision latency of \sys policies allows them to react to thermal
events before hardware protection mechanisms trigger.

\begin{figure}
\centering
\fbox{\parbox{\columnwidth}{\centering [Figure: Power-aware Throttling]}}
\caption{Power consumption and performance under thermal-aware throttling policies}\label{fig:power-throttle}
\end{figure}